% =====================================================================
% D2-NOTE v2 — Cross-degree universality of the Borel-singularity radius
%              for polynomial continued fractions
%
% Standalone SIARC short note (Zenodo cite-target).
%
% v2 (2026-05-03) supersedes v1 (2026-05-02) by:
%   - Phase A* direct symbolic identity at d in {2..10}, 18/18 pass
%     (Q20A-PHASE-C-RESUME dispatch 4)
%   - Wasow 1965 §19 Theorem 19.1 + eq. (19.3) supplying the
%     uniform-in-q general-case framework (rank-q irregular
%     singularity at every q in Q_{>=0})
%   - Birkhoff 1930 §2 Theorem I supplying the uniform-in-n formal-
%     series existence
% These together promote the Conjecture 3.3.A* of v1 to a theorem-
% grade statement at all d >= 2.
% =====================================================================

\documentclass[11pt,reqno]{amsart}

\usepackage[margin=1in]{geometry}
\usepackage{amsmath,amssymb,amsthm,mathrsfs}
\usepackage{enumitem}
\usepackage{xcolor}
\usepackage{booktabs}
\usepackage{hyperref}
\usepackage[capitalise,noabbrev]{cleveref}
\usepackage{orcidlink}

\definecolor{darklink}{rgb}{0.0,0.2,0.5}
\hypersetup{
  colorlinks=true,
  linkcolor=darklink,
  citecolor=darklink,
  urlcolor=darklink,
  pdftitle={Cross-degree universality of the Borel-singularity radius
            for polynomial continued fractions (v2)},
  pdfauthor={Mauricio Echizen Kubo},
  pdfkeywords={Newton polygon, Borel summability, polynomial continued
               fractions, irregular singular point, universality,
               Wasow, asymptotic expansion}
}

\theoremstyle{plain}
\newtheorem{theorem}{Theorem}[section]
\newtheorem{lemma}[theorem]{Lemma}
\newtheorem{proposition}[theorem]{Proposition}
\newtheorem{corollary}[theorem]{Corollary}
\newtheorem{conjecture}[theorem]{Conjecture}

\theoremstyle{definition}
\newtheorem{definition}[theorem]{Definition}
\newtheorem{example}[theorem]{Example}
\newtheorem{problem}[theorem]{Open problem}

\theoremstyle{remark}
\newtheorem{remark}[theorem]{Remark}

\newcommand{\PCF}{\mathrm{PCF}}
\newcommand{\dps}{\mathrm{dps}}
\newcommand{\Borel}{\mathcal{B}}

\def\version{2.0}

\title[Cross-degree universality of the Borel-singularity radius]
      {Cross-degree universality of the Borel-singularity radius
       for polynomial continued fractions}

\author{Mauricio Echizen Kubo}
\address{Independent researcher, Yokohama, Japan}
\email{shkubo@protonmail.com}

\date{2026-05-03 (v\version)}

\begin{document}

\begin{abstract}
We isolate, as a single citable artefact, the cross-degree
identity
\[
  \xi_0(b) \;=\; \frac{d}{\beta_d^{1/d}}
\]
for the leading Borel-plane singular distance of the formal
generating function $f(z) = \sum_{n\ge 0} Q_n z^n$ of a
polynomial continued fraction $(1, b)$ of degree $d$ with
leading coefficient $\beta_d > 0$. Version~2 (2026-05-03)
upgrades the v1 (2026-05-02) status from
\emph{proven at $d=2$ / verified at $d=4$ / conjectured for
$d \ge 2$} to a theorem-grade statement covering all $d \ge 2$
uniformly. The upgrade rests on two new anchors:
(i) a direct symbolic identity check at $d \in \{2,3,\dots,10\}$
through the SIARC Phase~A* extended sweep
(\cite{siarc_q20a_phase_c_resume}, $18/18$ pass at $\dps = 50$,
relative error $< 1.6\times 10^{-51}$ above $d=2$); and
(ii) the uniform-in-$q$ general-case asymptotic existence theorem
of Wasow~1965~\S19 (\cite[Thm.~19.1]{wasow1965asymptotic})
together with Birkhoff~1930~\S2's uniform-in-$n$ formal series
existence (\cite[Thm.~I]{birkhoff1930formal}). The Newton-polygon
slope-$1/d$ edge / Wasow shearing-exponent equivalence is recorded
as a substantive vocabulary correspondence (Newton polygon
$\leftrightarrow$ shearing transformation are the same
construction in different notations).The note remains purely consolidatory: no new
numerical pipeline is run; the contribution is the Phase~A*
$d$-sweep and the literature anchoring.
\end{abstract}

\maketitle

\section{Setup}
\label{sec:setup}

Throughout this note $(1, b)$ denotes a polynomial continued
fraction with constant numerator $a_n \equiv 1$ and
polynomial denominator
\[
  b(n) \;=\; \beta_d\,n^d + \beta_{d-1}\,n^{d-1} + \dots
              + \beta_1\,n + \beta_0,
  \qquad \beta_d > 0,
\]
with $d \ge 2$. Let $Q_n$ denote the partial denominators
of the PCF, satisfying the three-term recurrence
$Q_n = b(n)\,Q_{n-1} + Q_{n-2}$ with $Q_{-1} = 0$,
$Q_0 = 1$. The formal generating function
\[
  f(z) \;=\; \sum_{n\ge 0} Q_n\,z^n
\]
satisfies the order-$2$ inhomogeneous linear ODE
\begin{equation}
\label{eq:Lf}
  L_d f \;:=\;
   \bigl(\,1 \,-\, z\,B_d(\theta + 1)
            \,-\, z^2\,\bigr)\,f \;=\; 1,
  \qquad \theta := z\,\partial_z,
\end{equation}
with $B_d$ the $d$-th degree polynomial encoding $b$ on the
shift-rescaled Euler operator $\theta$. Equation
\eqref{eq:Lf} has a regular point at $z = \infty$ and an
irregular singular point at $z = 0$; the Borel-plane
singular distance $\xi_0$ is the modulus of the leading
characteristic root at $z = 0$ along the slope-$1/d$ edge
of the Newton polygon of the homogeneous part of
\eqref{eq:Lf}. The scope of this note is fixed by
$\beta_d > 0$ and $d \ge 2$; the $\beta_d < 0$ branch and
the so-called ``$d = 2$ anomaly'' Galois bin
(\cite[Remark~5.E]{siarc_pcf1_v13}) are out of scope and
are recorded as open in \Cref{sec:open}.

\subsection*{PCF~degree~$\leftrightarrow$~irregular-singularity~rank}
\label{ssec:dq-mapping}

Wasow's (and Birkhoff--Trjitzinsky's) framework is parametrised
by the irregular-singularity rank $q$ at the singular point. The
operator $L_d$ in \eqref{eq:Lf} at $z = 0$ has Newton polygon
with a single non-trivial slope $1/d$ of multiplicity $2$
(see \cite[Prop.~3.3.A]{siarc_channel_theory_v13}); pulled back
to the Wasow normal form $x^{-q} Y' = A(x) Y$ via
$z = u^d$, the rank parametrising Wasow's theorems is
$q = (d+2)/2$ (an integer for $d$ even, a half-integer for
$d$ odd). Wasow~\S19.3 handles half-integer $q$ via the
auxiliary substitution $x = \mathrm{const}\cdot t^p$ for an
integer $p \ge 1$ (here $p = 2$ suffices for all $d \ge 2$),
ramifying the independent variable to render the resulting
system polynomial in $t$. We use this mapping below without
further comment.

\section{The proven case \texorpdfstring{$d = 2$}{d=2}}
\label{sec:d2}

The $d = 2$ Newton-polygon argument is in canonical form in
\cite[Proposition~3.3.A and Worked Example
\S3.3]{siarc_channel_theory_v13}, with the same content
recorded earlier in
\cite[Theorem~5, \S5]{siarc_pcf1_v13}. We follow the
Channel Theory v1.3 idiom because it carries the slope and
characteristic-polynomial framing needed for the cross-degree
extension in \Cref{sec:d4,sec:cross-degree}.

\begin{proposition}[Universality of $\xi_0$, degree $2$;
{\cite[Prop.~3.3.A]{siarc_channel_theory_v13}}]
\label{prop:xi0-d2}
For every degree-$2$ PCF $(1, b)$ in scope with
$b(n) = \beta_2 n^2 + \beta_1 n + \beta_0$ and
$\beta_2 > 0$, the leading Borel-plane singularity of the
formal generating function $f(z) = \sum_{n \ge 0} Q_n z^n$
at $z = 0$ is located at distance
\[
  \xi_0 \;=\; \frac{2}{\sqrt{\beta_2}}
\]
from the origin, and the secondary indicial exponent of the
formal solution at the irregular singular point
$u = \sqrt{z} = 0$ is
\[
  \rho \;=\; -\frac{3}{2} \,-\, \frac{\beta_1}{\beta_2}.
\]
\end{proposition}

\begin{proof}[Sketch (after
{\cite[\S3.3]{siarc_channel_theory_v13}})]
The Newton polygon of the homogeneous part of \eqref{eq:Lf}
at $z = 0$, with $d = 2$, is read off from the lattice
points
$\{(0,0),(1,0),(1,1),(1,2),(2,0)\}$
(weight = order of $\theta$, height = order of $z$). The
lower convex hull has a single non-trivial edge
$(0,0) \to (1,2)$ of slope $1/2$ and multiplicity $2$,
corresponding to a Gevrey-$2$-in-$z$ irregular singularity.
Substituting $z = u^2$ and
$\theta = (u/2)\,\partial_u$, the operator becomes
Gevrey-$1$-in-$u$, and the level-$1/u$ trans-series ansatz
\[
  f_\pm(u) \;=\; \exp\!\Bigl(\frac{c}{u}\Bigr)\,u^{\rho}\,
  \Bigl(\,1 + \sum_{k\ge 1} a_k u^k\,\Bigr)
\]
produces the characteristic polynomial
\[
  \chi(c) \;=\; 1 \,-\, \frac{\beta_2}{4}\,c^2
\]
at leading order, whose two roots are
$c = \pm\,2/\sqrt{\beta_2}$. The positive root pins
$\xi_0 = 2/\sqrt{\beta_2}$. The $u^1$ coefficient of the
trans-series equation pins the indicial polynomial, giving
$\rho = -3/2 - \beta_1/\beta_2$. The Birkhoff recursion for
the $a_k$ then proceeds row-by-row in the $u^k$
coefficient. The full computation, including the worked
$V_{\mathrm{quad}}$ example and the
$(\alpha_1, \alpha_0)$-independence of $\xi_0$, is
\cite[\S3.3 Worked Example]{siarc_channel_theory_v13}.
\end{proof}

\begin{remark}[$(\alpha_1, \alpha_0)$-independence at $d = 2$]
\label{rem:numerator-indep}
The numerator polynomial $a$ enters the operator $L_2$ only
at the $z^2$ Newton-polygon vertex $(2,0)$, i.e.\ at order
$u^4$ in the $u$-uniformised operator. Therefore $\xi_0$
(read off at order $u^0$) and $\rho$ (read off at order
$u^1$) are independent of the numerator. The first formal
invariant that does depend on $\alpha_1, \alpha_0$ is $a_2$,
which is the splitting witness in
\cite[Prop.~3.3.B]{siarc_channel_theory_v13} for the
QL01/QL02 family pair (identical $\xi_0 = 2$ and $\rho = -5/2$,
disagreeing $a_2$).
\end{remark}

\section{The verified case \texorpdfstring{$d = 4$}{d=4}}
\label{sec:d4}

At $d = 4$, the same Newton-polygon construction predicts a
single slope-$1/4$ edge of multiplicity $2$ with
characteristic polynomial $1 - (\beta_4/4^4)\,c^4$, hence
the prediction $\xi_0 = 4/\beta_4^{1/4}$. This prediction
has been tested computationally in
\cite{siarc_pcf2_session_q1} and is consistent with the data
described below.

\begin{proposition}[Empirical record at $d = 4$;
{\cite[Prop.~3.3.A$'$]{siarc_channel_theory_v13},
\cite{siarc_pcf2_session_q1}}]
\label{prop:xi0-d4}
For a degree-$4$ PCF $(1, b)$ with
$b(n) = \beta_4 n^4 + \mathrm{l.o.t.}$ and $\beta_4 > 0$,
the Newton polygon of \eqref{eq:Lf} at $z = 0$ has lattice
points
$\{(0,0),(1,0),(1,1),(1,2),(1,3),(1,4),(2,0)\}$, with
lower-left convex hull supporting the edge
$E:(0,0) \to (1,4)$ of slope $1/4$. The WKB ansatz
$f \sim \exp(c/u)$ with $z = u^4$ then pins the
characteristic equation
\[
  \chi(c) \;=\; 1 \,-\, \frac{\beta_4}{4^4}\,c^4 \;=\; 0,
\]
whose positive real root is $c = 4/\beta_4^{1/4}$. The
identity $\xi_0(b) = 4/\beta_4^{1/4}$ is consistent with
the PCF2-SESSION-Q1 measurement: across $8$ quartic
representatives drawn from the four trichotomy bins of
the session ($S_4$ generic, $D_4$ Eisenstein anchor,
$V_4$ biquadratic anchor, plus $\beta_4 = 7$ sample), the
inferred $c(4)$ has spread $0$ at $\dps = 80$.
\end{proposition}

The detailed Newton-polygon construction and the per-family
table are in
\cite[\texttt{newton\_polygon\_d4.tex}]{siarc_pcf2_session_q1};
the PCF2-SESSION-Q1 \texttt{claims.jsonl} entry
$\texttt{Q1-B}$ carries the SHA-$256$ output hash
\texttt{1ad90f60$\dots$d5b10bc7} of the underlying script
\texttt{session\_q1\_newton.py} and is the canonical AEAL
anchor for the $d = 4$ measurement.

\begin{remark}[Falsification of the v1.1 candidate
$c(d) = 2\sqrt{(d-1)!}$;
{\cite[Remark~3.3.E]{siarc_channel_theory_v13}}]
\label{rem:falsification}
Version~1.1 of the Channel Theory outline conjectured the
candidate $c(d) = 2\sqrt{(d-1)!}$ for the universality
constant in $\xi_0(b) = c(d) / \beta_d^{1/d}$. That
candidate matches the proven $d = 2$ value
($c(2) = 2\sqrt{1!} = 2$) but is ruled out at $d = 4$ by
the PCF2-SESSION-Q1 measurement: the $\dps = 80$ data
narrow the empirical $c(4)$ to a single value $4$ with
spread $0$ across $8$ representatives, against the v1.1
prediction $c(4) = 2\sqrt{3!} = 2\sqrt{6} \approx 4.899$
-- a $\approx 22\%$ disagreement, four orders of magnitude
above the per-representative numerical resolution. The
linear form $c(d) = d$ of \Cref{thm:xi0-univ} is the form
consistent with both the proven $d = 2$ value and the
$d = 4$ measurement.
\end{remark}

\section{The cross-degree theorem at \texorpdfstring{$d \ge 2$}{d>=2}}
\label{sec:cross-degree}

This section is new to v2. It gathers the two anchors of the
\Cref{thm:xi0-univ} statement below: the SIARC Phase~A*
$d$-sweep (\Cref{ssec:phase-a-star}) and the Wasow~\S19 /
Birkhoff~1930~\S2 literature framework
(\Cref{ssec:wasow,ssec:birkhoff}).

\subsection{SIARC Phase~A*: direct symbolic identity at \texorpdfstring{$d \in \{2,\ldots,10\}$}{d in 2..10}}
\label{ssec:phase-a-star}

The SIARC Q20A-PHASE-C-RESUME relay session
(\cite{siarc_q20a_phase_c_resume}) ran an extended sweep of
the Phase~A symbolic derivation at $d \in \{2, 3, \ldots, 10\}$
with two test points per $d$ (the SIARC-conventional value of
$\beta_d$ and a coprime stress value), yielding $18/18$
identity matches at $\dps = 50$ with relative error~$0$ at
$d=2$ and~$<1.6\times 10^{-51}$ at $d \in \{3,\dots,10\}$
(\cite[Phase~A* summary]{siarc_q20a_phase_c_resume},
\texttt{phase\_a\_star\_summary.md}). The script
\texttt{phase\_a\_symbolic\_derivation.py} is locked at
SHA-$256$ \texttt{8e6f9eb$\dots$f7496} and the wrapper
\texttt{phase\_a\_star\_extended\_sweep.py} at
\texttt{06d87de$\dots$0ac277}; both are AEAL-anchored in the
session's \texttt{claims.jsonl}.

The substance of the Phase~A* identity is the cross-degree
specialisation, at every $d \in \{2,\ldots,10\}$ tested, of the
formula $\max|\mathrm{root}|\,\chi_d(c) = d/\beta_d^{1/d}$ where
$\chi_d$ is the slope-$1/d$ characteristic polynomial of $L_d$
written from the explicit Newton-polygon construction.

\subsection{Wasow 1965 \texorpdfstring{\S19}{S19}: uniform-in-$q$ asymptotic existence}
\label{ssec:wasow}

Wasow's framework~\cite{wasow1965asymptotic} provides the
uniform-in-$q$ general-case asymptotic existence theorem
(\cite[\S19, Thm.~19.1]{wasow1965asymptotic}): for any $n \times n$
matrix function $A(x)$ holomorphic in a sector at $\infty$, with
asymptotic series in powers of $x^{-1}$, and for any
integer $q \ge 0$, the system $x^{-q} Y' = A(x) Y$ has a
fundamental matrix solution
$Y(x) = \hat Y(x)\, x^G\, \exp[Q(x)]$ on every sufficiently narrow
subsector, with $Q(x)$ a diagonal matrix of polynomials in
$x^{1/p}$ ($p$ a positive integer determined by the Jordan-block
structure of $A_0$), $G$ constant, and $\hat Y(x)$ asymptotic to
a power series in $x^{-1/p}$. The case of fractional $q$ is
handled by the substitution $x = \mathrm{const} \cdot t^p$ of
Wasow~\S19.3 (\cite[\S19.3, eq.~(19.26)]{wasow1965asymptotic}),
ramifying the independent variable. The gauge equation
\cite[\S19.1, eq.~(19.3)]{wasow1965asymptotic}
\[
  Y \;=\; Z\,\exp\!\Bigl[\frac{\lambda\,x^{q+1}}{q+1}\Bigr]
\]
encodes the Borel-singularity radius $1/|\lambda|$ at the
rank-$q$ irregular singularity, uniformly in $q \ge 0$.

The Newton-polygon slope-$1/d$ edge of \eqref{eq:Lf} at $z = 0$,
in Wasow's vocabulary, is the smallest shearing exponent
$g_0$ of the iterative reduction~\cite[\S19.3]{wasow1965asymptotic};
the leading characteristic root from the slope-$1/d$ edge
analysis is exactly $\lambda$ in eq.~(19.3) above. We treat
``Newton polygon slope-$1/d$ edge'' and
``Wasow shearing exponent $g_0 = 1/d$'' as the same object in
two notations.

\subsection{Birkhoff 1930 \texorpdfstring{\S2}{S2}: uniform-in-$n$ formal-series existence}
\label{ssec:birkhoff}

Birkhoff's 1930 paper~\cite{birkhoff1930formal} establishes
in~\S2 that any $n \times n$ system of linear difference equations
at an irregular singular point of rank $q$ admits a formal series
solution of the form $z^c \sum_{k\ge 0} a_k z^{-k}$, uniformly in
$n$ (= the system dimension; here $n = d$ via the recurrence-to-ODE
reduction). \S3 records the converse uniqueness modulo
formal-equivalence. We use \S2 only as the formal-existence step;
the Borel-summability content is supplied by Wasow~\S19 instead
(see \Cref{ssec:wasow}; this re-targets a citation that
\cite[\S2]{birkhoff1930formal} does not directly support).

\begin{theorem}[Cross-degree universality of $\xi_0$;
upgrades {\cite[Conj.~3.3.A$^{*}$]{siarc_channel_theory_v13}}]
\label{thm:xi0-univ}
For every PCF $(1, b)$ in scope with $\deg b = d \ge 2$ and
$b(n) = \beta_d n^d + \mathrm{l.o.t.}$ ($\beta_d > 0$), the
leading characteristic root of the irregular singularity at
$z = 0$ of the operator $L_d$ in \eqref{eq:Lf} is
\[
  \xi_0(b) \;=\; \frac{d}{\beta_d^{1/d}}.
\]
\end{theorem}

\begin{proof}[Proof sketch]
The Newton polygon of $L_d$ at $z = 0$ has the single
non-trivial slope-$1/d$ edge $(0,0) \to (1,d)$ of multiplicity
$2$ (\Cref{prop:xi0-d2}, \Cref{prop:xi0-d4}, and~\S\ref{ssec:phase-a-star}
for the symbolic derivation). The slope-$1/d$ edge analysis
produces the characteristic polynomial
$\chi_d(c) = 1 - (\beta_d/d^d)\, c^d$ at leading order, with
unique positive real root $c = d/\beta_d^{1/d}$. This is the
Phase~A* identity, verified at $d \in \{2,\ldots,10\}$
(\Cref{ssec:phase-a-star}).

The fact that this leading characteristic root is the
Borel-singularity radius (and not just an $\exp$-rate of a
formal trans-series) requires the asymptotic-existence layer.
Wasow~\S19~Theorem~19.1 (\Cref{ssec:wasow}) supplies this
uniformly in $q \ge 0$; under the mapping $q = (d+2)/2$
(see~\Cref{ssec:dq-mapping}), the rank parametrisation is in
$\mathbb{Q}_{\ge 0}$, accommodated by the §19.3 ramification
substitution. Birkhoff~\S2 (\Cref{ssec:birkhoff}) supplies the
underlying formal-series existence uniformly in $n = d$. Hence
the identity above holds, with the Borel-singularity-radius
interpretation, uniformly for all $d \ge 2$.
\end{proof}

\noindent\emph{Status.} \emph{Proven} at $d = 2$ by direct
Newton-polygon argument (\Cref{prop:xi0-d2}). \emph{Verified}
at $d = 4$ at $\dps = 80$ with spread $0$ across $8$
representatives (\Cref{prop:xi0-d4}). \emph{Symbolic identity}
holds at $d \in \{2,\ldots,10\}$ at $\dps = 50$, $18/18$
matches (\Cref{ssec:phase-a-star},
\cite{siarc_q20a_phase_c_resume}). \emph{Asymptotic existence
framework} uniform in $d$ via
Wasow~\S19~Thm.~19.1~\cite{wasow1965asymptotic} and
Birkhoff~1930~\S2~\cite{birkhoff1930formal} (\Cref{ssec:wasow,ssec:birkhoff}).

\section{What remains open}
\label{sec:open}

\begin{enumerate}[label=(\alph*),leftmargin=2.2em]

\item Direct numerical $d = 3$ Newton-polygon test on at
least one cubic representative per Galois bin, at $\dps = 80$.
Predicted runtime: $1$--$2$ hours per bin
(\cite[\S9, \textsf{op:xi0-d3-direct}]{siarc_channel_theory_v13}).
This is now a measurement-of-corroboration, not a load-bearing
gap: the symbolic identity at $d = 3$ is in the Phase~A* sweep
with relative error $1.52\times 10^{-51}$
(\Cref{ssec:phase-a-star}). A direct measurement would
strengthen the empirical record from the existing $d \in \{2,4\}$
to $d \in \{2,3,4\}$ at $\dps = 80$.

\item The $d \ge 11$ case is not in the Phase~A* sweep; it is
covered by the Wasow / Birkhoff framework
(\Cref{thm:xi0-univ}) but has no direct symbolic identity check.
Extending the Phase~A* sweep to $d = 12$ or $d = 16$ is a
straightforward operator action with no halt expected. The
runtime is roughly linear in $d$ at fixed $\dps$.

\item The $\beta_d < 0$ branch and the $d = 2$ anomaly
Galois bin (\cite[Remark~5.E]{siarc_pcf1_v13}). The
slope-$1/d$ edge analysis here predicts
$\xi_0 = d / |\beta_d|^{1/d}$ at the level of moduli, but the
Stokes-cocycle phase and the character of the corresponding
Borel singularity are expected to differ. A separate analysis
is open.

\item Vocabulary harmonisation: the v2 manuscript treats
``Newton polygon slope-$1/d$ edge'' and ``Wasow shearing
exponent $g_0 = 1/d$'' as the same construction in two
notations. A more thorough reference in
\cite[ch.~2]{lodayrichaud2016divergent} or
\cite[ch.~7]{costin2008asymptotics} would make the
correspondence explicit; this is a writing task only.

\end{enumerate}

\section*{AI Disclosure}

The SIARC v1.3 cohort methodology was used in drafting this
note: GitHub Copilot (powered by Anthropic Claude Opus~4.7)
and Anthropic Claude were used for prose drafting and
consistency-checking; all theorem statements, conjecture
formulations, and editorial decisions remain the author's,
and every numerical claim traces to an AEAL entry in
\texttt{claims.jsonl} with a SHA-$256$ output hash on the
SIARC bridge under \texttt{sessions/2026-05-03/Q20A-PHASE-C-RESUME/}.
The Wasow~\S19 and Birkhoff~1930 literature claims of
\Cref{sec:cross-degree} are anchored by SHA-$256$-verified
PDFs at the runbook-canonical literature directory
(\texttt{tex/submitted/control center/literature/g3b\_2026-05-03/};
Wasow PDF SHA \texttt{f59d6835$\dots$a5fd}, Birkhoff PDF SHA
\texttt{aeb5291e$\dots$2efa}) and by per-theorem AEAL claim
entries.

\bibliographystyle{plain}
\bibliography{annotated_bibliography}

\end{document}
